\subsection{Hessenberg Reduction and QR Eigenvalue Methods}

\begin{definition} \label{def:upper-hessenberg-matrix} \hlt{Upper Hessenberg Matrix}\\
Let $A=[a_{ij}]\in\operatorname{Mat}(n,n,\mathbb{F})$. Then $A$ is \hlt{upper Hessenberg} if $a_{ij}=0$ whenever $i>j+1$.\\
Equivalently, all entries below the first subdiagonal are zero.
\end{definition}

Numerical linear algebra also replaces a matrix by a unitarily similar one with many forced zero entries. Upper Hessenberg form is the standard preliminary reduction before QR-type eigenvalue algorithms.

\begin{figure}[H]
  \centering
  \begin{tikzpicture}[font=\small]
  \node at (0,0) {$
  \begin{bmatrix}
  * & * & * & * & *\\
  * & * & * & * & *\\
  0 & * & * & * & *\\
  0 & 0 & * & * & *\\
  0 & 0 & 0 & * & *
  \end{bmatrix}
  $};
  \draw[red!70, thick] (-1.35,-0.45) -- (0.65,-1.15);
  \node[red!70!black, align=center] at (2.45,-1.0) {first\\subdiagonal};
  \draw[->, red!70, thick] (1.65,-0.95) -- (0.72,-0.92);

  \node[blue!70!black, align=center] at (0,1.75) {entries below this subdiagonal are zero};
  \draw[->, blue!70, thick] (-0.2,1.45) -- (-1.0,-0.85);
\end{tikzpicture}

  \caption{Upper Hessenberg pattern: red is the first subdiagonal; blue entries below it are forced zero.}
  \label{fig:upper-hessenberg-pattern}
\end{figure}

\begin{theorem} \label{thm:unitary-upper-hessenberg-reduction} \hlt{Unitary Upper Hessenberg Reduction}\\
Every $A\in\operatorname{Mat}(n,n,\mathbb{F})$ is unitarily similar to upper Hessenberg matrix with real nonnegative subdiagonal entries.
\begin{proof}
Set $A_0=A$. For $k=1,\ldots,n-2$, suppose columns $1,\ldots,k-1$ of $A_{k-1}$ already have zeros below the first subdiagonal. Let $\mathbf{w}$ be the tail of column $k$ in rows $k+1,\ldots,n$.\\
By Corollary \ref{cor:householder-unitary-annihilator}, choose a unitary $V_k$ with $V_k^*\mathbf{w}=\alpha\mathbf{e}_1$, $\alpha\geq0$.\\
Put $U_k=I_k\oplus V_k$ and $A_k=U_k^*A_{k-1}U_k$. Right multiplication by $U_k$ changes only columns $k+1,\ldots,n$, so columns $1,\ldots,k-1$ keep their existing Hessenberg zeros. Left multiplication by $U_k^*$ acts only on rows $k+1,\ldots,n$; those rows are already zero in columns $1,\ldots,k-1$, so the zeros remain. In column $k$, the tail becomes $V_k^*\mathbf{w}=\alpha\mathbf{e}_1$, so entries below row $k+1$ are zero and the new subdiagonal entry is real nonnegative.\\
After $n-2$ steps, $A_{n-2}=U^*AU$ is upper Hessenberg, where $U=U_1\cdots U_{n-2}$ is unitary.\\
\hlt{Structural Commentary}: Each trailing unitary creates one Hessenberg column and preserves earlier ones.\\
\hlt{Key Insight}: Kill one column tail at a time.
\end{proof}
\end{theorem}

\begin{perspective} \label{per:hessenberg-before-qr} \hlt{Why Hessenberg Form Matters}\\
Hessenberg form preserves spectral data by unitary similarity, making QR-type eigenvalue algorithms cheaper.
\end{perspective}

\begin{method} \label{meth:qr-eigenvalue-algorithm-preview} \hlt{QR Eigenvalue Algorithm Preview}\\
For a square matrix $A_0$, the basic QR eigenvalue iteration is:
\begin{enumerate}[label=(\roman*)]
\setlength{\itemsep}{0pt}
\item In practice, first reduce $A_0$ to Hessenberg form by Theorem \ref{thm:unitary-upper-hessenberg-reduction}.
\item Factor $A_k=Q_kR_k$, with $Q_k$ unitary and $R_k$ upper triangular.
\item Set $A_{k+1}=R_kQ_k=Q_k^*A_kQ_k$, using shifts to speed convergence.
\end{enumerate}
Each step is a unitary similarity, so the spectrum is preserved.
\end{method}

\begin{method} \label{meth:hessenberg-qr-step-plane-rotations} \hlt{Hessenberg QR Step by Plane Rotations}\\
Let $H\in\operatorname{Mat}(n,n,\mathbb{C})$ be upper Hessenberg and let $\mu\in\mathbb{C}$.
\begin{enumerate}[label=(\roman*)]
\setlength{\itemsep}{0pt}
\item Set $B_0=H-\mu I_n$.
\item For $j=1,\ldots,n-1$, choose a unitary plane rotation $G_j$ acting on coordinates $j,j+1$,\\
so that $(G_j^*B_{j-1})_{j+1,j}=0$ and $B_j=G_j^*B_{j-1}$.
\item Put $Q=G_1\cdots G_{n-1}$ and $R=B_{n-1}$; then $H-\mu I_n=QR$.
\item Define the shifted QR update $H_+=RQ+\mu I_n=Q^*HQ$
\end{enumerate}
Adjacent rotations implement the Hessenberg step by local bulge chasing.
\end{method}
